Hidden states are the fundamental information bottleneck in transformer language models, yet the relationship between decoder complexity and storage capacity remains poorly understood. We present the first systematic comparison of five encoding schemes spanning the full decoder-complexity spectrum---from zero-decoder bit-packing to full frozen-model decoding---applied to a single 768-dimensional \gptwo hidden state. Raw ASCII bit-packing stores the most information (24,576 bits at fp32) but requires no learned structure. LLM vocabulary packing stores 1,536 tokens losslessly. At the other extreme, a single learned \memvec vector decoded through the full frozen \gptwo model recovers $\sim$37 of 50 random tokens (578 bits), while 32 \memvec vectors achieve 99.4\% accuracy on 500 random tokens (7,762 bits) with peak efficiency of 0.46 bits per learnable parameter. We identify a sharp 3-dimension-per-token threshold for unembedding-based decoding, show that superposition encoding fails entirely for dense token IDs despite its success for sparse binary features, and find that a single frozen transformer layer adds minimal capacity beyond position embeddings alone. Our results reveal a clear decoder--capacity tradeoff: simpler encodings store more raw bits, but model-based decoders leverage learned structure to recover meaningful tokens at lower bit cost, with natural text amplifying this advantage dramatically.
